% !TEX program = xelatex
\documentclass[a4paper,11pt]{article}
\usepackage[UTF8]{ctex}
\setCJKmainfont{SimSun}[BoldFont=SimHei,ItalicFont=KaiTi]
\setCJKsansfont{SimHei}

% 页面布局与页边距
\usepackage[top=2.5cm,bottom=2.5cm,left=2.5cm,right=2.5cm]{geometry}

% 数学公式与符号
\usepackage{amsmath,amssymb,bm}
\usepackage{mathrsfs}

% 插图与绘图
\usepackage{graphicx}
\usepackage{tikz}
\usetikzlibrary{arrows.meta,calc}

% 表格与版面美化
\usepackage{booktabs}
\usepackage{xcolor}
\usepackage{fancyhdr}
\usepackage{tcolorbox}
\tcbuselibrary{skins,breakable}

% 定义颜色
\definecolor{primary}{HTML}{0D47A1}
\definecolor{secondary}{HTML}{1565C0}
\definecolor{accent}{HTML}{E65100}
\definecolor{bglight}{HTML}{F7FAFF}
\definecolor{border}{HTML}{D2ECD3}

% 页眉页脚设置
\pagestyle{fancy}
\fancyhf{}
\fancyhead[L]{\fontsize{9}{11}\selectfont 可压缩流体力学专题解析}
\fancyhead[R]{\fontsize{9}{11}\selectfont 激波与膨胀波流动}
\fancyfoot[C]{\thepage}
\renewcommand{\headrulewidth}{0.4pt}

% 自定义环境
\newtcolorbox{problembox}[1]{
  colback=bglight,
  colframe=primary,
  fonttitle=\bfseries\large,
  title=#1,
  boxrule=1pt,
  arc=4pt,
  breakable
}

\newtcolorbox{keypointbox}[1]{
  colback=orange!5,
  colframe=accent,
  fonttitle=\bfseries,
  title=#1,
  boxrule=0.8pt,
  arc=3pt,
  breakable
}

\begin{document}

\title{\textbf{\Large 可压缩流动与气体动力学经典真题解析：\\ 激波与膨胀波的组合流动}}
\author{流体力学学习笔记}
\date{\today}
\maketitle

\thispagestyle{plain}

\begin{problembox}{原题重现}
在拍摄纹影图像有限大小的圆形窗口中，发现壁面 $A$ 处有极小划痕，显示出一道 $30^\circ$ 的波；拐角 $B$ 处有 $50^\circ$ 的波。如图 \ref{fig:problem_diagram} 所示。
\begin{enumerate}
    \item 求壁面偏折角 $\theta$；
    \item 已知 $CD$ 上的静压为 $p_3 = 1.0 \times 10^5 \text{ N/m}^2$，求来流高压大储气罐内的压力 $p_0$。
\end{enumerate}
\textbf{已知条件：} 气体为无粘常比热完全气体（比热比 $\gamma = 1.4$），$AB$、$CD$ 段壁面与水平来流速度平行。
\end{problembox}

\begin{figure}[htbp]
\centering
\begin{tikzpicture}[scale=1.3, >=Stealth]
  % 定义坐标
  \coordinate (O) at (0,0);
  \coordinate (A) at (1.5,0);
  \coordinate (B) at (4.5,0);
  \coordinate (C) at (6.5,0.6548);
  \coordinate (D) at (8.5,0.6548);
  
  % 绘制壁面
  \draw[line width=1.5pt, line cap=rect] (O) -- (A) -- (B) -- (C) -- (D);
  
  % 壁面下方的阴影线
  \foreach \x in {0.0, 0.2, ..., 4.4} {
    \draw[thin, gray] (\x, 0) -- (\x-0.15, -0.15);
  }
  \foreach \t in {0.0, 0.1, ..., 1.0} {
    \coordinate (P) at ($(B)!\t!(C)$);
    \draw[thin, gray] (P) -- ($(P) + (-0.1,-0.15)$);
  }
  \foreach \x in {6.5, 6.7, ..., 8.5} {
    \draw[thin, gray] (\x, 0.6548) -- (\x-0.15, 0.5048);
  }
  
  % 辅助虚线
  \draw[dashed, gray] (B) -- ++(2.5, 0);
  \draw[dashed, gray] (C) -- ++(-1.5, 0);
  \draw[dashed, gray] (C) -- ++(2.0, 0);
  
  % 壁面偏折角 theta 标注
  \draw[draw=black, thin, ->] (4.5+1.2, 0) arc (0:18.13:1.2);
  \node[right] at (4.5+1.25, 0.2) {$\theta$};
  
  % A处划痕
  \fill[red] (A) circle (1.5pt) node[below=2pt] {A (划痕)};
  
  % A处的马赫波 (30度)
  \draw[orange, thick, ->] (A) -- ++(30:2.8) node[above right] {马赫波 ($30^\circ$)};
  \draw[thin] (A) ++(0.8, 0) arc (0:30:0.8);
  \node[right] at ($(A)+(30:0.8)$) {\fontsize{8}{10}\selectfont $30^\circ$};
  
  % B处的斜激波 (50度)
  \draw[red, thick, ->] (B) -- ++(50:2.8) node[above right] {斜激波 ($50^\circ$)};
  \draw[thin] (B) ++(0.8, 0) arc (0:50:0.8);
  \node[above right] at ($(B)+(0.8,0.1)$) {\fontsize{8}{10}\selectfont $50^\circ$};
  
  % 拐角标注
  \node[below=2pt] at (B) {B};
  \node[above=2pt] at (C) {C};
  
  % C处的膨胀扇形区
  \draw[blue, thin] (C) -- ++(68.01:2.0);
  \draw[blue, thin] (C) -- ++(58.6:2.0);
  \draw[blue, thin] (C) -- ++(49.2:2.0);
  \draw[blue, thin] (C) -- ++(39.8:2.0);
  \draw[blue, thin] (C) -- ++(31.17:2.0);
  \node[blue, right] at ($(C)+(39.8:2.0)$) {膨胀波};
  
  % 区域标注
  \node at (0.8, 0.8) {\textbf{区域 \uppercase\expandafter{\romannumeral1}}};
  \node at (0.8, 0.5) {$M_1 = 2$};
  \node at (0.8, 0.2) {$p_1, p_{01}$};
  
  \node at (5.0, 1.2) {\textbf{区域 \uppercase\expandafter{\romannumeral2}}};
  \node at (5.0, 0.9) {$M_2 \approx 1.31$};
  \node at (5.0, 0.6) {$p_2, p_{02}$};
  
  \node at (8.0, 1.8) {\textbf{区域 \uppercase\expandafter{\romannumeral3}}};
  \node at (8.0, 1.5) {$M_3 \approx 1.93$};
  \node at (8.0, 1.2) {$p_3, p_{03}$};
  
  % 纹影圆形窗口
  \draw[dashed, black!60, thick] (3.0, 0.5) circle (1.8);
  \node[black!60] at (3.0, 2.5) {拍摄区域（圆形窗口）};

\end{tikzpicture}
\caption{流动结构与物理区域划分示意图}
\label{fig:problem_diagram}
\end{figure}

\section{核心物理模型与基础公式}

本题属于典型的\textbf{超声速气体动力学流动问题}，主要涉及以下三个核心物理过程：

\begin{keypointbox}{物理过程分解与公式链}
\begin{enumerate}
    \item \textbf{微弱扰动与马赫波 (A 处)：}
    壁面上的极小划痕仅引起微弱扰动，因此在超声速流场中会产生一条马赫线。马赫线与来流方向的夹角为马赫角 $\mu$，其与来流马赫数 $M$ 的关系为：
    \begin{equation}
        \sin \mu = \frac{1}{M}
    \end{equation}
    
    \item \textbf{压缩拐角与斜激波 (B 处)：}
    超声速流通过内折角（压缩拐角）时，会产生斜激波以偏转气流。偏折角 $\theta$、斜激波角 $\beta$ 和激波前马赫数 $M_1$ 满足经典的 $\theta-\beta-M$ 关系式：
    \begin{equation}
        \tan \theta = 2 \cot \beta \frac{M_1^2 \sin^2 \beta - 1}{M_1^2 (\gamma + \cos 2\beta) + 2}
    \end{equation}
    激波前后的法向马赫数分别为 $M_{n1} = M_1 \sin \beta$ 和 $M_{n2}$，其中 $M_{n2}$ 的关系式为：
    \begin{equation}
        M_{n2} = \sqrt{\frac{1 + \frac{\gamma-1}{2} M_{n1}^2}{\gamma M_{n1}^2 - \frac{\gamma-1}{2}}}
    \end{equation}
    激波后的实际马赫数 $M_2$ 为：
    \begin{equation}
        M_2 = \frac{M_{n2}}{\sin(\beta - \theta)}
    \end{equation}
    激波前后的总压（滞止压强）之比 $\frac{p_{02}}{p_{01}}$ 可以根据滞止压强的定义以及激波前后的静压比直接推导。因为：
    \begin{equation}
        p_0 = p \left( 1 + \frac{\gamma-1}{2} M^2 \right)^{\frac{\gamma}{\gamma-1}}
    \end{equation}
    因此，激波后的总压 $p_{02}$ 与激波前的总压 $p_{01}$ 之比为：
    \begin{equation}
        \frac{p_{02}}{p_{01}} = \frac{p_2}{p_1} \times \left( \frac{1 + \frac{\gamma-1}{2} M_2^2}{1 + \frac{\gamma-1}{2} M_1^2} \right)^{\frac{\gamma}{\gamma-1}}
    \end{equation}
    其中，激波前后的静压比 $\frac{p_2}{p_1}$ 可以由正激波静压比公式（以法向马赫数 $M_{n1} = M_1 \sin \beta$ 为自变量）决定：
    \begin{equation}
        \frac{p_2}{p_1} = 1 + \frac{2\gamma}{\gamma+1} (M_{n1}^2 - 1)
    \end{equation}
    对于空气（常比热比 $\gamma = 1.4$），上述总压比公式可简写为极具物理意义且易于记忆的式子：
    \begin{equation}
        \frac{p_{02}}{p_{01}} = \frac{p_2}{p_1} \times \left( \frac{1 + 0.2 M_2^2}{1 + 0.2 M_1^2} \right)^{3.5}
    \end{equation}
    
    \item \textbf{膨胀拐角与 Prandtl-Meyer 膨胀波 (C 处)：}
    超声速流通过外折角（膨胀拐角）时，会通过 Prandtl-Meyer 膨胀扇形区进行无旋等熵膨胀。膨胀角 $\Delta \theta$ 与马赫数的关系为：
    \begin{equation}
        \nu(M_3) = \nu(M_2) + \Delta \theta
    \end{equation}
    其中 $\nu(M)$ 为 Prandtl-Meyer 函数（对于 $\gamma = 1.4$ 的空气）：
    \begin{equation}
        \nu(M) = \sqrt{6} \arctan \left( \sqrt{\frac{M^2-1}{6}} \right) - \arctan \left( \sqrt{M^2-1} \right)
    \end{equation}
    由于膨胀过程为等熵过程，气流的滞止压力保持不变：
    \begin{equation}
        p_{03} = p_{02}
    \end{equation}
    气流在任意马赫数 $M$ 下的静压 $p$ 与滞止压强 $p_0$ 满足等熵关系：
    \begin{equation}
        \frac{p_0}{p} = \left( 1 + \frac{\gamma-1}{2} M^2 \right)^{\frac{\gamma}{\gamma-1}}
    \end{equation}
\end{enumerate}
\end{keypointbox}

\section{详细求解步骤}

\subsection{第一步：求解来流马赫数 $M_1$}

在壁面 $A$ 处，极小划痕产生的微弱波属于马赫波。因为壁面 $AB$ 平行于来流方向，所以马赫波与水平面夹角即为马赫角：
\[
\mu_1 = 30^\circ
\]
由此可以求得来流马赫数 $M_1$：
\begin{equation}
    M_1 = \frac{1}{\sin \mu_1} = \frac{1}{\sin 30^\circ} = 2.0
\end{equation}

\subsection{第二步：求解壁面偏折角 $\theta$ （第一问）}

在拐角 $B$ 处，壁面向上折角 $\theta$，气流受到压缩产生斜激波。斜激波与水平来流的夹角为激波角：
\[
\beta = 50^\circ
\]
将 $M_1 = 2.0$, $\beta = 50^\circ$, $\gamma = 1.4$ 代入 $\theta-\beta-M$ 关系式中：
\begin{align*}
    \tan \theta &= 2 \cot 50^\circ \frac{2.0^2 \sin^2 50^\circ - 1}{2.0^2 (1.4 + \cos 100^\circ) + 2}
\end{align*}
计算各个分项：
\begin{itemize}
    \item $\cot 50^\circ \approx 0.8391$
    \item $\sin 50^\circ \approx 0.7660 \implies \sin^2 50^\circ \approx 0.5868$
    \item $\cos 100^\circ \approx -0.1736$
\end{itemize}
代入分子与分母：
\begin{align*}
    \text{分子} &= 2 \times 0.8391 \times (4.0 \times 0.5868 - 1) = 1.6782 \times 1.3473 \approx 2.2610 \\
    \text{分母} &= 4.0 \times (1.4 - 0.1736) + 2 = 4.0 \times 1.2264 + 2 = 6.9056 \\
    \tan \theta &= \frac{2.2610}{6.9056} \approx 0.3274
\end{align*}
由此解得壁面偏折角 $\theta$ 为：
\begin{equation}
    \theta = \arctan(0.3274) \approx 18.13^\circ
\end{equation}

\subsection{第三步：求解斜激波后的流动状态（区域 \uppercase\expandafter{\romannumeral2}）}

为了向下游反推，首先需要求得区域 \uppercase\expandafter{\romannumeral2} 的法向马赫数 $M_{n1}$：
\[
M_{n1} = M_1 \sin \beta = 2.0 \sin 50^\circ \approx 1.5321
\]
激波后法向马赫数 $M_{n2}$：
\[
M_{n2} = \sqrt{\frac{1 + 0.2 M_{n1}^2}{1.4 M_{n1}^2 - 0.2}} = \sqrt{\frac{1 + 0.2 \times 2.3473}{1.4 \times 2.3473 - 0.2}} = \sqrt{\frac{1.4695}{3.0862}} \approx 0.6900
\]
则激波后区域 \uppercase\expandafter{\romannumeral2} 的实际马赫数 $M_2$ 为：
\begin{equation}
    M_2 = \frac{M_{n2}}{\sin(\beta - \theta)} = \frac{0.6900}{\sin(50^\circ - 18.13^\circ)} = \frac{0.6900}{\sin(31.87^\circ)} \approx \frac{0.6900}{0.5280} \approx 1.3069
\end{equation}

首先计算激波前后的静压比为：
\begin{equation}
    \frac{p_2}{p_1} = 1 + \frac{2 \times 1.4}{1.4 + 1} (M_{n1}^2 - 1) = 1 + \frac{7}{6} (1.5321^2 - 1) \approx 2.5718
\end{equation}
接着根据滞止压强的定义比值计算激波前后的总压比为：
\begin{align}
    \frac{p_{02}}{p_{01}} &= \frac{p_2}{p_1} \times \left( \frac{1 + 0.2 M_2^2}{1 + 0.2 M_1^2} \right)^{3.5} \nonumber \\
    &= 2.5718 \times \left( \frac{1 + 0.2 \times 1.3069^2}{1 + 0.2 \times 2.0^2} \right)^{3.5} \nonumber \\
    &= 2.5718 \times \left( \frac{1 + 0.2 \times 1.7080}{1 + 0.2 \times 4.0} \right)^{3.5} \nonumber \\
    &= 2.5718 \times \left( \frac{1.3416}{1.8000} \right)^{3.5} \approx 2.5718 \times 0.7453^{3.5} \nonumber \\
    &\approx 2.5718 \times 0.3574 \approx 0.9193
\end{align}
（斜激波导致了约 $8.07\%$ 的总压损失。）

\begin{tcolorbox}[colback=yellow!5,colframe=orange,title=学习与考场提示：如何快速获取此总压比值？]
在实际计算中，该方法不仅极具物理意义（直接源于滞止压强定义），而且计算量明显小于记忆复杂的正激波总压公式。
\begin{itemize}
    \item \textbf{推荐方法：查《正激波表》（Normal Shock Table）}。在超声速空气动力学考试中，也可以直接通过查正激波表快速获得此值。以激波的法向马赫数分量 $M_{n1} = 1.5321$ 查正激波表：
    \[
    \text{查表得：} \quad M_{n1} = 1.53 \implies \frac{p_{02}}{p_{01}} \approx 0.9200; \quad M_{n1} = 1.54 \implies \frac{p_{02}}{p_{01}} \approx 0.9166
    \]
    对其进行简单的线性插值：
    \[
    \frac{p_{02}}{p_{01}} \approx 0.9200 - 0.21 \times (0.9200 - 0.9166) = 0.9193
    \]
    这与解析公式计算得出的 $0.9193$ 完全一致。
\end{itemize}
\end{tcolorbox}

\subsection{第四步：求解膨胀波后的流动状态（区域 \uppercase\expandafter{\romannumeral3}）}

在拐角 $C$ 处，壁面水平偏折角从上倾 $\theta$ 回归到水平 $CD$，气流向下偏折了：
\[
\Delta \theta = \theta = 18.13^\circ \approx 0.3164 \text{ rad}
\]
首先计算区域 \uppercase\expandafter{\romannumeral2} 的 Prandtl-Meyer 函数值 $\nu(M_2)$：
\begin{align*}
    \nu(M_2) &= \sqrt{6} \arctan \left( \sqrt{\frac{1.3069^2-1}{6}} \right) - \arctan \left( \sqrt{1.3069^2-1} \right) \\
    &= \sqrt{6} \arctan \left( \sqrt{\frac{0.7080}{6}} \right) - \arctan \left( \sqrt{0.7080} \right) \\
    &= \sqrt{6} \arctan(0.3435) - \arctan(0.8414) \\
    &\approx 2.4495 \times 0.3308 \text{ rad} - 0.6994 \text{ rad} \\
    &\approx 0.8103 \text{ rad} - 0.6994 \text{ rad} = 0.1109 \text{ rad} \approx 6.359^\circ
\end{align*}
通过 Prandtl-Meyer 膨胀，区域 \uppercase\expandafter{\romannumeral3} 的函数值为：
\begin{equation}
    \nu(M_3) = \nu(M_2) + \Delta \theta = 6.359^\circ + 18.13^\circ = 24.489^\circ \approx 0.4274 \text{ rad}
\end{equation}
根据 $\nu(M_3) = 24.489^\circ$，反查 Prandtl-Meyer 函数表或通过数值迭代求解 $M_3$：
\begin{equation}
    M_3 \approx 1.9320
\end{equation}

\subsection{第五步：反推来流滞止压强 $p_0$ （第二问）}

由于 $C$ 处的膨胀流动是等熵的，区域 \uppercase\expandafter{\romannumeral2} 到区域 \uppercase\expandafter{\romannumeral3} 的总压守恒：
\[
p_{03} = p_{02}
\]
已知 $CD$ 上的静压 $p_3 = 1.0 \times 10^5 \text{ N/m}^2$，根据等熵关系求区域 \uppercase\expandafter{\romannumeral3} 的总压 $p_{03}$：
\begin{align}
    p_{03} &= p_3 \left( 1 + \frac{\gamma-1}{2} M_3^2 \right)^{\frac{\gamma}{\gamma-1}} \nonumber \\
    &= 1.0 \times 10^5 \times \left( 1 + 0.2 \times 1.9320^2 \right)^{3.5} \nonumber \\
    &= 1.0 \times 10^5 \times (1 + 0.7465)^{3.5} \nonumber \\
    &= 1.0 \times 10^5 \times 7.0411 = 7.0411 \times 10^5 \text{ N/m}^2
\end{align}
因为等熵膨胀，故区域 \uppercase\expandafter{\romannumeral2} 的总压为：
\[
p_{02} = p_{03} = 7.0411 \times 10^5 \text{ N/m}^2
\]
由来流大储气罐的压力即来流滞止压力 $p_{01}$：
\begin{equation}
    p_{01} = \frac{p_{02}}{0.9193} = \frac{7.0411 \times 10^5}{0.9193} \approx 7.659 \times 10^5 \text{ N/m}^2
\end{equation}

因此，来流高压大储气罐内的压力为：
\begin{equation*}
    \bm{p_{01} \approx 7.66 \times 10^5 \text{ N/m}^2}
\end{equation*}

\section{方法学总结与知识点拓展}

\begin{enumerate}
    \item \textbf{为什么这不像“常见”的流体力学题？}
    普通流体力学课（如土木、水利、市政方向）通常只讲到\textbf{不可压缩流}（液体和低速气体，马赫数 $M < 0.3$）。这类课程的核心是伯努利方程、管路水头损失、粘性阻力等。
    本题是典型的\textbf{气体动力学 (Gas Dynamics) / 可压缩流体力学}问题，核心是气流的“波关系”。
    
    \item \textbf{马赫波与划痕的妙用：}
    纹影法 (Schlieren photography) 是利用光在不同密度流体中的折射率差异来显影波纹的技术。壁面上的微小划痕（极小扰动）在超声速气流中天然地产生一条马赫线，成了测量来流马赫数 $M_1$ 的“无源测量计”。
    
    \item \textbf{总压损失的实质：}
    等熵膨胀（C处）是不损失总压的，但是激波（B处）是强烈的非等熵过程。超声速气流经过激波时，由于强烈的粘性与热传导，机械能耗散，使得总压降低。因此在反推来流总压时，必须乘上激波总压恢复系数的倒数。
\end{enumerate}

\end{document}
